%! Composition of Functions — Unit 5, Lesson 5.6 — Student Packet Cover Sheet
\documentclass[10pt]{article}
\usepackage{algebra2-article}
\usepackage{algebra2-boxes}
\usepackage{ltablex}
\keepXColumns

\begin{document}

% Full-bleed forest banner
\begin{tikzpicture}[remember picture, overlay]
  \fill[forest] (current page.north west)
    rectangle ([yshift=-0.9in]current page.north east);
\end{tikzpicture}
\vspace{-0.8in}

\begin{center}
  {\color{white}\textbf{\LARGE Algebra 2}}\\[4pt]
  {\color{forestmid}\large\textbf{Unit 5 \quad Radical Functions}}\\[6pt]
  {\color{white}\textbf{\large Lesson 5.6 \quad Composition of Functions}}
\end{center}

\vspace{0.22in}
\namedateperiod
\vspace{0.18in}

\begin{learningtargetbox}
\small
\begin{enumerate}[leftmargin=1.5em, itemsep=5pt]
  \item \ldots{} read the notation $(f\circ g)(x)=f(g(x))$ correctly and evaluate a composition \textbf{numerically} (from values, from a table) and \textbf{graphically} (from two graphs), working \textbf{inside out}.
  \item \ldots{} find $(f\circ g)(x)$ \textbf{algebraically} by substituting one whole function in for the other's input --- and show that $(f\circ g)(x)$ and $(g\circ f)(x)$ are usually \emph{different functions}.
  \item \ldots{} determine the \textbf{domain of a composition} by checking both stages, and explain why the simplified formula can hide a restriction.
  \item \ldots{} take a function apart --- write $h(x)$ as $f(g(x))$ --- and recognize that every transformation from Lesson 5.4 was a composition all along.
\end{enumerate}
\end{learningtargetbox}

\vspace{0.14in}

\begin{tocbox}
\small
\renewcommand{\arraystretch}{1.5}
\begin{tabularx}{\linewidth}{@{} c l X r @{}}
\rowcolor{forestbg}
\textbf{\#} & \textbf{Component} & \textbf{Description} & \textbf{Score} \\
\midrule
1 & Warm-Up        & Spiral review: function notation, substituting a whole expression, and why order changes an answer & \blank{1.2cm} \\
2 & Guided Notes   & The notation, composing numerically / graphically / algebraically, order, domain, and taking a function apart & \blank{1.2cm} \\
3 & Group Activity & \emph{Wired in Series} --- composing, comparing the two orders, and finding the pair that undoes itself, by tier & \blank{1.2cm} \\
4 & Exit Ticket    & Quick check: both orders, one radical composition with its domain, plus an SOL-style item & \blank{1.2cm} \\
5 & Homework       & Values, a table, algebra with domains, two graphs, an error to diagnose, and the view from the mast & \blank{1.2cm} \\
\midrule
  & \multicolumn{2}{r}{\textbf{Total}} & \blank{1.2cm} \\
\end{tabularx}
\end{tocbox}

\vspace{0.10in}

\begin{remindbox}
\small Until today, every function you met worked alone. Today you wire two of them together: the
output of one becomes the input of the next. That is a \textbf{composition}, written
$(f\circ g)(x)=f(g(x))$, and the first thing to get right is the \emph{order} --- the function
standing closest to $x$ runs \textbf{first}, so you work \textbf{inside out}. Order is not a
formality here. Take $20\%$ off a jacket and then hand over a $\$10$ coupon and you pay one price;
hand over the coupon first and you pay a different one. In symbols, $(f\circ g)(x)$ and
$(g\circ f)(x)$ are usually two \emph{different functions} --- and this unit gives the cleanest
example there is: with $f(x)=\sqrt{x}$ and $g(x)=x-4$, one order gives $\sqrt{x-4}$ and the other
gives $\sqrt{x}-4$. You already know both of those graphs from Lesson 5.4, and they do not even have
the same domain. That is the second idea of the day: a composition happens in \textbf{two stages},
so an input has to survive \emph{both} of them. Watch what that does to
$\left(\sqrt{x}\,\right)^{2}$ --- the formula simplifies all the way down to $x$, and yet you still
cannot put $-9$ into it, because stage one refused the input before stage two ever saw it. Hold on
to that. In Lesson 5.7 it turns out to be the whole reason $y=x^{2}$ needs a restricted domain
before it can have a square-root inverse.
\end{remindbox}

\end{document}
